\section{Discussion}
\label{sec:discussion}

% \subsection{On Baselines for Comparison}
% While we compare \appname against a strong LLM-based baseline (FlatLLM) in both defect modeling and impact analysis, we acknowledge that broader comparisons with existing tools are limited.
% This is primarily due to the lack of directly comparable systems tailored for the same problem setting.
% Existing PR analysis tools (e.g., commit summarizers, vulnerability scanners, or static detectors like Bandit, Semgrep) either focus on general-purpose bug summarization without downstream reasoning, or target traditional languages like Java with coarse-grained dependency analysis.
% To the best of our knowledge, no existing tool jointly models defect semantics and performs client-side impact analysis over DL libraries.
% We thus position FlatLLM as a strong representative of the LLM-first, non-agentic strategy that developers  naturally adopt.




% \subsection{On Client Dataset Size and Dynamic Validation}
% Our downstream impact analysis is conducted on 57 PR–client and 6 commit–client pairs. 
% While this may appear limited in scale, each sample requires substantial effort to prepare and evaluate.
% For every (defect, client) pair, we manually locate relevant open-source projects and validate trigger condition satisfaction via static and dynamic analysis.
% Dynamic reproduction, in particular, involves downloading models, configuring environments, and crafting test inputs that align with defect semantics—requiring hours per case in many instances.
% Despite the moderate size, our results reveal meaningful patterns and support actionable insights, including identifying two real-world clients silently affected by internal library defects.
% We believe this level of depth, while costly, provides strong evidence of \appname's practical value and correctness.


\subsection{Generalizability, Adoption, and Extension}
\textbf{Generalizability.}
While our evaluation focuses on two popular DL libraries—Transformers and \textit{Megatron}, the modular architecture of \appname is broadly applicable to other Python-based libraries that expose public APIs and have structured commits or PRs.
We observe that \appname performs consistently across a range of defect types (e.g., configuration misuse, semantic errors, performance issues), suggesting good generalization. Extending the system to non-Python libraries (e.g., TensorFlow C++) would require adapting AST parsing and static analysis routines.

\noindent \textbf{Developer Adoption.}
Case studies in \textit{Megatron} show that users may prefer to locally backport patches instead of upgrading the entire library version, due to compatibility or dependency constraints. 
In such scenarios, \appname’s ability to surface concrete trigger conditions and risk factors can assist developers in writing minimal fixes. 
This highlights its potential utility as a defect-awareness tool in long-lived or production-grade codebases.

\noindent\textbf{Extension Directions.}
One promising extension is to integrate dynamic tracing or execution logs to better infer runtime-specific defects (e.g., those triggered by hardware mismatch or environment settings).
Another direction is model-agnostic tuning, where different LLMs (e.g., GPT-4, DeepSeek, Claude) are compared in terms of consistency, robustness, and hallucination rates under the same agent roles.
We also plan to explore IDE integration to offer actionable defect awareness during development or CI workflows.


\subsection{Threats to Validity}
\textbf{Threats to Internal Validity.} Some of our annotations—such as defect presence or client impact—rely on expert judgment, though we observe substantial inter-annotator agreement.
Certain trigger conditions may also depend on runtime paths not visible in static code.
Although \appname performs consistency checks and retry rounds, occasional hallucinations from LLMs cannot be fully ruled out.
\textbf{Threats to External Validity.} Our evaluation is limited to two actively maintained DL libraries.
Libraries with less structured development workflows or sparse metadata may degrade the effectiveness of agent-based reasoning.
In addition, the current system is designed for Python-based ecosystems; applying it to other languages requires further adaptation.
